\documentclass[12pt]{article}
\usepackage[utf8]{inputenc}
\usepackage[T1]{fontenc}
\usepackage{hyperref}
\title{Puelche Language Resources for K-12 Curriculum Development}
\author{Compiled by Claude Sonnet 5.5 (claude-sonnet-5-5) for Andr\'as Kornai}
\date{}
\begin{document}
\maketitle

\section{Introduction}
Puelche (ISO 639-3 \texttt{pue}, Glottocode puel1244), also called Gününa Küne or Günün a Yajüch, is a dormant language of Argentina, usually treated as an isolate; some scholars have proposed links to the Chon languages or to the language of the Querandí, on thin evidence. Sources give different dates for its last fluent speakers (1934 or 1960). There are no native speakers, so any programme is heritage and revitalization work based on historical documentation.

\section{Primary Educational Resources}
None found. No school materials or primers for Puelche turned up in searches.

\section{Linguistic and Reference Materials}
The search results mention these historical and modern works by name, but I did not find their texts online, so no links are given for them: Hunziker's 1864 Genakenn vocabulary and phrase list (Vocabulario y fraseario Genakenn); Barbará's 1879 Manual o Vocabulario de la Lengua Pampa; and Casamiquela's 1983 Nociones de gramática del gününa küne. These should be consulted in libraries or archives.
\begin{itemize}
\item Gününa Küne language, Wikipedia (English): overview, classification, sources. \url{https://en.wikipedia.org/wiki/G%C3%BCn%C3%BCna_K%C3%BCne_language}
\item Het peoples, Wikipedia (English): background on the peoples of the region. \url{https://en.wikipedia.org/wiki/Het_peoples}
\item Günün-Chonan Languages, University of Edinburgh page. \url{https://chico-sc.ppls.ed.ac.uk/?p=425}
\item WALS language page for Gününa Küne. \url{https://wals.info/languoid/lect/wals_code_gku}
\item Endangered Languages Project, bibliography for this language. \url{https://endangeredlanguages.com/lang/2899/bibliography?hl=fr}
\end{itemize}

\section{Religious and Community Resources}
None found.

\section{Digital and Audio Resources}
None found. No audio exists, since the language was last spoken decades ago.

\section{Curriculum Development Context}
A heritage curriculum could draw on the historical word lists and Casamiquela's grammar, with the history and geography of the Puelche (Gününa Küne) people as cultural content, ideally with descendant communities. The documentation is thin, so such a programme would teach history and a small amount of reconstructed vocabulary rather than fluent language use.

\end{document}